\documentclass[aspectratio=169,11pt]{beamer}

\usepackage{booktabs}
\usepackage{array}
\usepackage{siunitx}
\usepackage{xurl}

\usetheme{Madrid}
\usecolortheme{default}
\setbeamertemplate{navigation symbols}{}
\setbeamertemplate{footline}[frame number]

\newcommand{\movetag}[2]{\textbf{Move #1:} #2}
\newcommand{\steptag}[2]{\textbf{Step #1:} #2}
\newcommand{\claim}[1]{\textcolor{blue!65!black}{#1}}
\newcommand{\gaptext}[1]{\textcolor{red!65!black}{#1}}
\newcommand{\solutiontext}[1]{\textcolor{green!45!black}{#1}}
\newcommand{\src}[1]{\textsuperscript{\scriptsize[\texttt{\detokenize{#1}}]}}
\newcommand{\key}[1]{\texttt{\detokenize{#1}}}
\newcolumntype{L}[1]{>{\raggedright\arraybackslash}p{#1}}

\title[MICE Introduction Map]{MICE Structure Map for the Introduction}
\subtitle{From the current situation to the thesis surrogate workflow}
\author{Linan Jiang}
\institute{Aalto University}
\date{7 April 2026}

\begin{document}

\begin{frame}
  \titlepage
\end{frame}

\begin{frame}{MICE Model: Four Moves in an Engineering Thesis Introduction}
\small
\begin{tabular}{L{0.18\linewidth}L{0.34\linewidth}L{0.38\linewidth}}
\toprule
Move & Function & Question answered in this thesis \\
\midrule
Move 1 & Current situation & Why is pilot spray in multi-fuel marine engines worth studying? \\
Move 2 & Problems & What is missing from experiments, CFD, empirical models, and image measurements? \\
Move 3 & New approach & Why is learning a surrogate from Mie videos a reasonable route? \\
Move 4 & Your solution & What does this thesis do, how does it do it, and what is its scope? \\
\bottomrule
\end{tabular}

\vspace{0.7em}
\footnotesize
The MICE model emphasizes a general-to-specific and problem-solution order. This deck uses it to audit the argument chain in the introduction.
\end{frame}

\begin{frame}{Overall Logic of the Introduction}
\small
\begin{enumerate}
  \item \claim{Industry background}: shipping decarbonization and alternative fuels drive multi-fuel marine-engine development.
  \item \claim{Technical background}: many multi-fuel engines still need pilot fuel for reliable ignition.
  \item \gaptext{Engineering risk}: pilot diesel spray development may lead to impingement, lubricant dilution, and deteriorated combustion.
  \item \gaptext{Method gap}: engine tests and CFD are costly, while image measurements alone do not form a fast predictive screening model.
  \item \solutiontext{New route}: build reduced-order penetration targets from high-speed Mie imaging and train an uncertainty-aware surrogate.
\end{enumerate}
\end{frame}

\section{Move 1}

\begin{frame}{Move 1-1: Making a Centrality Claim}
\small
\movetag{1}{Current situation}

\vspace{0.5em}
\steptag{1}{Making a centrality claim}

\vspace{0.8em}
\begin{block}{Thesis argument}
International shipping faces simultaneous pressure to reduce greenhouse-gas emissions, diversify fuels, and maintain operational reliability. The IMO 2023 strategy and industry orderbook trends make alternative-fuel propulsion a practical engineering issue \src{imo2023ghgstrategy; classnk2025alternativefuels}.
\end{block}

\begin{itemize}
  \item Function: places the work in the wider context of marine energy transition.
  \item Keywords: importance, relevance, application.
  \item Introduction location: Research relevance, paragraph 1.
\end{itemize}
\end{frame}

\begin{frame}{Move 1-2: Making Topic Generalizations}
\small
\movetag{1}{Current situation}

\vspace{0.5em}
\steptag{2}{Making topic generalizations}

\vspace{0.8em}
\begin{block}{Thesis argument}
Dual-fuel and multi-fuel marine engines require not only alternative-fuel supply but also reliable ignition. Low-pressure gas, methanol, and ammonia concepts commonly still require a small quantity of pilot fuel \src{wingd_xdf_manual_2021; wingd_lng_faq_2025; classnk_methanol_2024; wingd_ammonia_faq_2025}.
\end{block}

\begin{itemize}
  \item Function: explains the general role of pilot fuel in the engine system.
  \item Narrowing step: pilot diesel spray is not only an ignition aid; its spatial development also creates engineering risk.
  \item Introduction location: Research relevance, paragraphs 2--3.
\end{itemize}
\end{frame}

\begin{frame}{Move 1-3: Reviewing Previous Research and Solutions}
\small
\movetag{1}{Current situation}

\vspace{0.5em}
\steptag{3}{Reviewing previous research and solutions}

\vspace{0.8em}
\begin{block}{Thesis argument}
Spray behavior and impingement risk are commonly assessed through engine tests, CFD and spray-combustion simulations, and optical spray experiments. Diesel-spray research already has a foundation in high-speed imaging, empirical correlations, phenomenological models, and CFD \src{Linne2013SprayImaging; HiroyasuArai1990; NaberSiebers1996; liu2020sprayreview}.
\end{block}

\begin{itemize}
  \item Function: acknowledges existing work instead of claiming that diesel spray is unstudied.
  \item Introduction location: Current engineering practice and Existing solutions.
\end{itemize}
\end{frame}

\section{Move 2}

\begin{frame}{Move 2-1: Indicating Weaknesses}
\small
\movetag{2}{Problems}

\vspace{0.5em}
\steptag{1}{Indicating weaknesses}

\vspace{0.8em}
\begin{itemize}
  \item Engine tests are closest to real operation, but they are expensive, slow, and hard to use for isolating one design factor.
  \item High-fidelity CFD is information-rich, but it requires complex models, many assumptions, and substantial computational resources.
  \item One reacting diesel-spray LES example required about 48,000 CPU core-hours and up to 20 days of wall-clock time for only \(\SI{2}{ms}\) after SOI \src{CurtisNiemeyerSung2017}.
  \item Empirical and phenomenological models are compact and interpretable, but less robust for short-pulse, strongly transient, and multiple-injection cases.
  \item Optical boundaries are method-dependent experimental definitions, not hidden physical truth fields \src{ecn_liquid_penetration; ecn_liquid_penetration_accessory}.
\end{itemize}
\end{frame}

\begin{frame}{Move 2-2: Identifying a Gap}
\small
\movetag{2}{Problems}

\vspace{0.5em}
\steptag{2}{Identifying a gap}

\vspace{0.8em}
\begin{block}{Thesis gap statement}
What is missing is a fast and repeatable method for predicting the time-evolving non-reacting liquid envelope of pilot diesel spray, so that later impingement-risk analysis has a defensible geometric input.
\end{block}

\begin{itemize}
  \item The gap is not that diesel spray has never been studied.
  \item The missing element is an intermediate-layer predictor between optical experiments and high-cost validation.
  \item The aim is not to replace CFD or engine tests, but to narrow the candidate space before using them.
\end{itemize}
\end{frame}

\section{Move 3}

\begin{frame}{Move 3-1: Introducing a Potential Solution}
\small
\movetag{3}{Proposing a new approach}

\vspace{0.5em}
\steptag{1}{Introducing a potential solution}

\vspace{0.8em}
\begin{block}{Thesis route}
Learn a condition-parametric surrogate directly from large-scale Mie-scattering optical spray data, turning measured penetration trajectories into a fast predictive model.
\end{block}

\begin{itemize}
  \item Mie videos already capture liquid-envelope behavior.
  \item Reduced-order fitting converts raw observations into stable learning targets.
  \item The surrogate sits between raw measurement and physics-based simulation.
\end{itemize}
\end{frame}

\begin{frame}{Move 3-2: Justifying the Approach}
\small
\movetag{3}{Proposing a new approach}

\vspace{0.5em}
\steptag{2}{Justifying the approach}

\vspace{0.8em}
\begin{itemize}
  \item \solutiontext{Experimental fidelity}: learning targets come from actual Mie images and penetration trajectories.
  \item \solutiontext{Screening speed}: avoids sending every candidate condition directly to CFD or engine testing.
  \item \solutiontext{Data extensibility}: the condition-to-trajectory-parameter mapping can be updated as the dataset grows.
  \item \solutiontext{Interpretability and auditability}: calibration, boundary definition, trajectory fitting, and uncertainty regression all leave intermediate artifacts.
  \item \solutiontext{Consistency with scientific imaging practice}: Fiji/ImageJ, scikit-image, OpenPIV, and PIVlab emphasize scriptable, inspectable, reusable pipelines \src{Schindelin2012Fiji; VanDerWalt2014ScikitImage; OpenPIV; ThielickeSonntag2021PIVlab}.
\end{itemize}
\end{frame}

\section{Move 4}

\begin{frame}{Move 4-1: Research Aims}
\small
\movetag{4}{Your solution}

\vspace{0.5em}
\steptag{1}{Research aims}

\vspace{0.8em}
\begin{block}{Thesis aim}
Develop a data-driven non-reacting liquid-envelope prediction model for pilot diesel spray. The model predicts spray-envelope and penetration trends from injection-condition and time information, together with uncertainty estimates suitable for engineering interpretation.
\end{block}

\begin{itemize}
  \item Output: mean penetration trend plus uncertainty.
  \item Use: early-stage impingement-risk screening in marine dual-fuel engine development.
\end{itemize}
\end{frame}

\begin{frame}{Move 4-2: Methods}
\small
\movetag{4}{Your solution}

\vspace{0.5em}
\steptag{2}{Methods}

\vspace{0.8em}
\begin{enumerate}
  \item High-speed Mie imaging of diesel spray in a room-temperature pressurized-air chamber.
  \item Manual calibration of nozzle center and pixel scale.
  \item CUDA-accelerated batch processing to extract plume-resolved temporal observables and boundary geometry.
  \item Reduced-order fitting to construct penetration targets and handle right-censoring and temporal alignment.
  \item Multi-stage learning: Stage-1 MSE, Stage-2 heteroscedastic NLL, and teacher-guided refinement.
  \item Coupling predicted mean and uncertainty to simplified geometric risk screening.
\end{enumerate}
\end{frame}

\begin{frame}{Move 4-3: Thesis Scope}
\small
\movetag{4}{Your solution}

\vspace{0.5em}
\steptag{3}{Thesis scope}

\vspace{0.8em}
\begin{itemize}
  \item In scope: controlled pressurized-chamber, room-temperature, non-reacting, non-evaporating pilot diesel spray.
  \item In scope: macroscopic liquid-envelope geometry and time evolution as observed in the optical setup.
  \item Out of scope: combustion chemistry, evaporation, dual-fuel ignition coupling, real hot-wall interaction, and real-engine wall-film truth.
  \item Model boundary: mainly interpolation within the covered data domain; clear OOD extrapolation is not guaranteed at the same level.
\end{itemize}
\end{frame}

\begin{frame}{Move 4-4: Main Outcome}
\small
\movetag{4}{Your solution}

\vspace{0.5em}
\steptag{4}{Describing the main outcome}

\vspace{0.8em}
\begin{enumerate}
  \item A CUDA-accelerated workflow from terabyte-scale Mie videos to plume-resolved observables.
  \item Reduced-order fitting for trajectory alignment, robust extrapolation, and mitigation of right-censoring.
  \item A multi-stage surrogate combining MSE, heteroscedastic NLL, and teacher-guided refinement.
  \item A probabilistic screening prototype that decouples penetration prediction from downstream geometric collision evaluation.
\end{enumerate}
\end{frame}

\begin{frame}{Move 4-5: Thesis Structure}
\small
\movetag{4}{Your solution}

\vspace{0.5em}
\steptag{5}{Overview of thesis structure}

\vspace{0.8em}
\begin{tabular}{L{0.27\linewidth}L{0.60\linewidth}}
\toprule
Part & Function \\
\midrule
Literature review and technical rationale & Situates the work in diesel spray imaging, optical diagnostics, and surrogate modeling. \\
Materials and methods & Describes the experimental data, image processing, feature extraction, trajectory fitting, and model training. \\
Results & Reports archived processing assets and model evidence. \\
Conclusions & Summarizes the outcome, limitations, and remaining development needs. \\
\bottomrule
\end{tabular}
\end{frame}

\begin{frame}{Move/Step Checklist}
\scriptsize
\begin{tabular}{L{0.18\linewidth}L{0.29\linewidth}L{0.41\linewidth}}
\toprule
MICE position & Introduction subsection & Core argument \\
\midrule
Move 1-1 & Research relevance & Shipping decarbonization and alternative fuels make the topic important. \\
Move 1-2 & Research relevance & Multi-fuel engines rely on pilot fuel, and spray spatial development creates risk. \\
Move 1-3 & Current practice / Existing solutions & Experiments, CFD, optical diagnostics, and empirical models already exist. \\
Move 2-1 & Existing solutions & Existing methods have weaknesses in cost, transience, interpretability, or truth definition. \\
Move 2-2 & Problem statement & A fast intermediate-layer predictor is missing between optical experiments and high-cost validation. \\
Move 3-1 & End of problem statement & Learn a surrogate from Mie data. \\
Move 3-2 & Why this workflow is necessary & Reduced targets, auditability, uncertainty, and high-throughput processing justify the route. \\
Move 4-1--5 & Aim / Methods / Scope / Contributions / Structure & The thesis solution, boundary, outcomes, and organization. \\
\bottomrule
\end{tabular}
\end{frame}

\begin{frame}{Source Keys}
\small
This deck uses citation keys from the thesis bibliography as lightweight source markers. The formal references remain in the thesis PDF bibliography.

\vspace{0.7em}
\begin{tabular}{L{0.34\linewidth}L{0.54\linewidth}}
\toprule
Topic & Source keys \\
\midrule
Shipping and alternative fuels & \key{imo2023ghgstrategy}; \key{classnk2025alternativefuels} \\
Pilot-fuel background & \key{wingd_xdf_manual_2021}; \key{wingd_lng_faq_2025}; \key{classnk_methanol_2024}; \key{wingd_ammonia_faq_2025} \\
Spray experiments and models & \key{Linne2013SprayImaging}; \key{HiroyasuArai1990}; \key{NaberSiebers1996}; \key{liu2020sprayreview} \\
CFD cost & \key{CurtisNiemeyerSung2017} \\
Optical boundary definition & \key{ecn_liquid_penetration}; \key{ecn_liquid_penetration_accessory} \\
Image-processing ecosystems & \key{Schindelin2012Fiji}; \key{VanDerWalt2014ScikitImage}; \key{OpenPIV}; \key{ThielickeSonntag2021PIVlab} \\
\bottomrule
\end{tabular}
\end{frame}

\end{document}
